\section{Commuting contributions within ordered gaps}

\begin{theorem}[Span of the commuting-gap forms]
    \label{thm:polynomial_commuting_gap_span}
    \lean{Polynomial.coeff_pow_C_add_C_mul_X_mem_span_commuting_gaps}
    \leanok
    Let $\mathcal R$ be a semiring equipped with a module structure over a
    semiring $\mathcal K$. Let $A,B,T\in\mathcal R$ satisfy $AB=BA$, and let
    $m,k\ge0$ be integers. For nonnegative exponents with
    $\sum_{j=0}^{k}(a_j+b_j)=m$, put
    \begin{align*}
        G(a,b)=A^{a_0}B^{b_0}T\cdots T A^{a_k}B^{b_k}.
    \end{align*}
    Then the coefficient $[X^k](A+B+TX)^{m+k}$ belongs to the
    $\mathcal K$-span of this fixed family of ordered products.
    The variable $X$ is central. The family includes zero exponents and
    need not be linearly independent.
    This is the commuting-gap step used in
    \cite[Lemma~4.2]{Arad2013AreaLaw}.
\end{theorem}
\begin{proof}
    \leanok
    \uses{thm:polynomial_linear_factor_coefficient}
    Let $\mathcal C_{m,k}$ be the additive submonoid generated by the displayed
    products. Left multiplication by $A$ increments $a_0$; left multiplication
    by $B$ increments $b_0$, using $BA^{a_0}=A^{a_0}B$. Left multiplication by
    $T$ inserts a first gap with both exponents zero. Consequently,
    \begin{align*}
        A\mathcal C_{m,k}+B\mathcal C_{m,k}
            &\subseteq\mathcal C_{m+1,k}, \\
        T\mathcal C_{m,k}&\subseteq\mathcal C_{m,k+1}.
    \end{align*}
    Write $F_{m,k}=[X^k](A+B+TX)^{m+k}$.
    Theorem~\ref{thm:polynomial_linear_factor_coefficient} gives
    \begin{align*}
        F_{m+1,k+1}&=AF_{m,k+1}+BF_{m,k+1}+TF_{m+1,k}, \\
        F_{m+1,0}&=AF_{m,0}+BF_{m,0},\qquad F_{0,k}=T^k.
    \end{align*}
    Induction on $m$ and then on $k$ proves $F_{m,k}\in\mathcal C_{m,k}$.
    Every $\mathcal K$-submodule containing the generators contains this
    additive submonoid, which proves the assertion.
    Repeated contributions are added: for example, $F_{2,0}=A^2+2AB+B^2$.
\end{proof}

\begin{theorem}[Number of indexed commuting-gap forms]
    \label{thm:polynomial_commuting_gap_count}
    \lean{Finset.card_piAntidiag_univ_fin_prod_bool}
    \leanok
    For nonnegative integers $m,k$, the number of tuples
    $(a_0,b_0,\ldots,a_k,b_k)$ of nonnegative integers with total $m$ is
    \begin{align*}
        \binom{m+2k+1}{2k+1}.
    \end{align*}
    Thus for $\ell=m+k$ the count is $\binom{\ell+k+1}{2k+1}$.
    This counts exponent tuples, even when their associated operators agree.
    The additional $+1$ corrects the printed count in
    \cite[Lemma~4.2]{Arad2013AreaLaw}; see
    \cite{gap:aklv13_interpolation_and_word_count}.
\end{theorem}
\begin{proof}
    \leanok
    There are $2(k+1)$ nonnegative coordinates with total $m$.
    Stars and bars gives $\binom{m+2k+1}{m}$, equal to the displayed count
    by the symmetry of binomial coefficients.
\end{proof}


\begin{theorem}[Operator Schmidt rank of a commuting-gap coefficient]
    \label{thm:polynomial_commuting_gap_rank}
    \lean{Matrix.operatorSchmidtRank_coeff_pow_C_add_C_mul_X_le}
    \leanok
    Let $A,B,T$ be operators on a finite-dimensional complex bipartite space.
    Suppose $AB=BA$, $\operatorname{OSR}(A),\operatorname{OSR}(B)\le1$, and
    $\operatorname{OSR}(T)\le r$ for an integer $r\ge0$. For integers $m,k\ge0$,
    \begin{align*}
        \operatorname{OSR}\!\left([X^k](A+B+TX)^{m+k}\right)
            \le\binom{m+2k+1}{2k+1}r^k.
    \end{align*}
    Zero-dimensional factors, zero exponents and $r=0$ are included, with
    the usual convention $0^0=1$.
\end{theorem}
\begin{proof}
    \leanok
    \uses{thm:polynomial_commuting_gap_span, thm:polynomial_commuting_gap_count,
        thm:operator_schmidt_rank_algebra}
    Every power of $A$ or $B$ has rank at most one, including the zeroth
    power. The product inequality therefore bounds each ordered gap form by $r^k$.
    The coefficient belongs to the span of the fixed finite family of these
    forms. Choose its scalar coefficients in that family. The scalar and
    finite-sum rank bounds give at most the sum of the ranks of the indexed
    forms, hence the displayed count times $r^k$.
\end{proof}


\begin{theorem}[Coefficient extraction in one tensor factor]
    \label{thm:local_polynomial_coefficient_rank}
    \lean{Matrix.operatorSchmidtRank_coeff_coeff_pow_left_add_right_le}
    \leanok
    Let $\mathcal H_L,\mathcal H_R$ be finite-dimensional complex spaces,
    let $A$ act on $\mathcal H_L$, let $B(Y)$ be a polynomial with operator
    coefficients on $\mathcal H_R$, and let $T$ act on
    $\mathcal H_L\otimes\mathcal H_R$ with $\operatorname{OSR}(T)\le r$.
    In the iterated polynomial algebra with central variables $Y$ and then $X$, put
    \begin{align*}
        F(X,Y)&=A\otimes\mathbb 1+\mathbb 1\otimes B(Y)+TX, \\
        K_{m,j,s}&=[Y^s]\bigl([X^j]F(X,Y)^{m+j}\bigr).
    \end{align*}
    For natural numbers $m,j,s,r$,
    \begin{align*}
        \operatorname{OSR}(K_{m,j,s})
            \le\binom{m+2j+1}{2j+1}r^j.
    \end{align*}
    Zero-dimensional factors and zero exponents are included, with $0^0=1$.
    This is supporting algebra for the coefficient restrictions in
    \cite[Lemma~4.2]{Arad2013AreaLaw} and
    \cite[Proposition~5.7]{Arad2016RigorousRG}.
\end{theorem}
\begin{proof}
    \leanok
    \uses{thm:polynomial_commuting_gap_span,thm:polynomial_commuting_gap_count,
        lem:operator_schmidt_block_span_bound,thm:operator_schmidt_rank_algebra}
    The left operator $A\otimes\mathbb 1$ commutes with
    $\mathbb 1\otimes B(Y)$. Apply the commuting-gap expansion over the
    coefficient algebra of operator polynomials in $Y$, retaining complex
    scalar coefficients. For an exponent tuple $(a_0,b_0,\ldots,a_j,b_j)$
    with total $m$, choose a shortest decomposition
    $T=\sum_{t=1}^{r_T}U_t\otimes V_t$, where $r_T=\operatorname{OSR}(T)$.
    Expanding only the $j$ occurrences of $T$ gives terms
    $L_{d,t}\otimes R_{d,t}(Y)$, where
    \begin{align*}
        L_{d,t}&=\left(\prod_{u=0}^{j-1}A^{a_u}U_{t_{u+1}}\right)A^{a_j}, \\
        R_{d,t}(Y)&=\left(\prod_{u=0}^{j-1}B(Y)^{b_u}V_{t_{u+1}}\right)B(Y)^{b_j}.
    \end{align*}
    The products are taken in increasing order of $u$.
    The left factors are independent of $Y$, so taking the coefficient of
    $Y^s$ leaves at most $r_T^j$ product terms for each exponent tuple.
    No factors are reordered across an occurrence of $T$.
    The scalar and finite-sum rank bounds, followed by the exponent-tuple
    count, give the claimed inequality. This argument bounds the actual
    extracted coefficient; it does not infer that coefficient extraction
    preserves arbitrary operator Schmidt rank bounds.
\end{proof}


\begin{theorem}[Interchanging the order of two marked coefficients]
    \label{thm:polynomial_two_marker_power}
    \lean{Polynomial.coeff_coeff_pow_two_markers_swap}
    \leanok
    Let $\mathcal R$ be a semiring and let $D,U,V\in\mathcal R$.
    With central polynomial variables, form
    \begin{align*}
        F_X(X)&=D+VY+UX\in\mathcal R[Y][X], \\
        F_Y(Y)&=D+UX+VY\in\mathcal R[X][Y].
    \end{align*}
    For natural numbers $k,i,j$,
    \begin{align*}
        [Y^j]\bigl([X^i]F_X(X)^k\bigr)
            =[X^i]\bigl([Y^j]F_Y(Y)^k\bigr).
    \end{align*}
    The elements $D,U,V$ may coincide or satisfy algebraic relations.
    The identity compares the same marked counts in the two nesting orders;
    it requires no commutation of $D,U,V$.
    It is supporting algebra for the simultaneous coefficient restrictions in
    \cite[Proposition~5.7]{Arad2016RigorousRG}.
\end{theorem}
\begin{proof}
    \leanok
    \uses{thm:polynomial_linear_factor_coefficient}
    Let $c_{k,i,j}$ denote either coefficient, with value zero when a lower
    index is negative. At $k=0$ its value is one if $i=j=0$ and zero otherwise.
    Applying Theorem~\ref{thm:polynomial_linear_factor_coefficient} in each
    polynomial variable gives
    \begin{align*}
        c_{k+1,i,j}
            =Dc_{k,i,j}+Uc_{k,i-1,j}+Vc_{k,i,j-1}.
    \end{align*}
    Both nesting orders satisfy this same recurrence and initial condition.
    Induction on $k$ proves the identity, including the cases $i=0$ or $j=0$.
    The proof changes only the order of addition, retaining the order of every
    product. For example, the mixed coefficient at $k=2$ is $UV+VU$.
\end{proof}


\begin{theorem}[A tripartite double-frequency rank bound]
    \label{thm:tripartite_polynomial_coefficient_rank}
    \lean{Matrix.operatorSchmidtRank_middle_coeff_coeff_pow_le}
    \leanok
    \uses{def:bipartite_operator_schmidt_rank}
    Let $\mathcal H_L,\mathcal H_M,\mathcal H_R$ be finite-dimensional complex
    spaces. Let $A_L,A_M,A_R$ act on the corresponding single factors, and
    let $U$ act on $\mathcal H_L\otimes\mathcal H_M$ and $V$ on
    $\mathcal H_M\otimes\mathcal H_R$. On the three-factor space, put
    \begin{align*}
        D&=A_L\otimes\mathbb 1_M\otimes\mathbb 1_R
            +\mathbb 1_L\otimes A_M\otimes\mathbb 1_R
            +\mathbb 1_L\otimes\mathbb 1_M\otimes A_R, \\
        \widetilde U&=U\otimes\mathbb 1_R, \\
        \widetilde V&=\mathbb 1_L\otimes V.
    \end{align*}
    With central variables $Y$ and then $X$, define
    \begin{align*}
        F(X,Y)&=D+\widetilde VY+\widetilde UX, \\
        K_{k,i,j}&=[Y^j]\bigl([X^i]F(X,Y)^k\bigr).
    \end{align*}
    Let $r_U,r_V,k,i,j$ be natural numbers satisfying $i\le k$, $j\le k$,
    $\operatorname{OSR}_{L\mid M}(U)\le r_U$ and
    $\operatorname{OSR}_{M\mid R}(V)\le r_V$. Then
    \begin{align*}
        \operatorname{OSR}_{M\mid LR}(K_{k,i,j})
            \le\binom{k+i+1}{2i+1}\binom{k+j+1}{2j+1}r_U^i r_V^j.
    \end{align*}
    The spaces may have dimension zero, with $0^0=1$.
    The binomial factors use the corrected exponent-tuple count in
    \cite{gap:aklv13_interpolation_and_word_count}.
    This is supporting algebra for the simultaneous boundary estimates in
    \cite[Proposition~5.7]{Arad2016RigorousRG}.
\end{theorem}
\begin{proof}
    \leanok
    \uses{def:bipartite_operator_schmidt_rank,
        thm:local_polynomial_coefficient_rank,thm:polynomial_two_marker_power,
        thm:operator_schmidt_middle_rank}
    Choose shortest product decompositions of $U$ and $V$. Tensoring each
    summand with the spectator identity gives a decomposition of the same
    length at its corresponding boundary. Reversing the two factors also
    preserves the length of a product decomposition. Thus the lifted crossing
    ranks are at most $r_U$ and $r_V$ in the respective boundary coordinates.

    At the left boundary the ordinary right polynomial is
    $B_L(Y)=A_M\otimes\mathbb 1_R+\mathbb 1_M\otimes A_R+VY$.
    Apply Theorem~\ref{thm:local_polynomial_coefficient_rank} with the
    constant left operator $A_L$, crossing operator $\widetilde U$,
    outer frequency $i$, remaining degree $k-i$, and inner frequency $j$.
    This gives
    \begin{align*}
        \operatorname{OSR}_{L\mid MR}(K_{k,i,j})
            \le\binom{k+i+1}{2i+1}r_U^i.
    \end{align*}
    At the reflected right boundary the ordinary opposite-factor polynomial
    is $B_R(X)=A_L\otimes\mathbb 1_M+\mathbb 1_L\otimes A_M+UX$.
    Theorem~\ref{thm:polynomial_two_marker_power} identifies the same
    $K_{k,i,j}$ when $V$ is counted by the outer variable.
    Applying Theorem~\ref{thm:local_polynomial_coefficient_rank} with
    outer frequency $j$, remaining degree $k-j$, and inner frequency $i$
    therefore gives
    \begin{align*}
        \operatorname{OSR}_{LM\mid R}(K_{k,i,j})
            \le\binom{k+j+1}{2j+1}r_V^j.
    \end{align*}
    Both inequalities concern the same operator. Their product bounds its
    middle-versus-exterior rank by
    Theorem~\ref{thm:operator_schmidt_middle_rank}.
\end{proof}

\begin{theorem}[Weighted frequency sums across the middle cut]
    \label{thm:weighted_full_frequency_middle_rank}
    \lean{Matrix.operatorSchmidtRank_middle_sum_smul_full_counts_le}
    \leanok
    \uses{def:bipartite_operator_schmidt_rank}
    Let $I$ be finite and let $H_a$ act on
    $\mathcal H_L\otimes\mathcal H_M\otimes\mathcal H_R$ for $a\in I$,
    where the three complex spaces are finite-dimensional.
    Fix distinct $u,v\in I$ and put $O=I\setminus\{u,v\}$.
    Suppose that $A_{L,a},A_{M,a},A_{R,a}$ act on the corresponding
    single factors and satisfy, for every $a\in O$,
    \begin{align}
        H_a & =A_{L,a}\otimes\mathbb 1_M\otimes\mathbb 1_R
        +\mathbb 1_L\otimes A_{M,a}\otimes\mathbb 1_R
        +\mathbb 1_L\otimes\mathbb 1_M\otimes A_{R,a}.\notag
    \end{align}
    Let $U$ act on $\mathcal H_L\otimes\mathcal H_M$ and $V$ on
    $\mathcal H_M\otimes\mathcal H_R$, with
    $H_u=U\otimes\mathbb 1_R$ and $H_v=\mathbb 1_L\otimes V$.
    Let $r_U,r_V,k,i,j$ be natural numbers such that
    $\operatorname{OSR}_{L\mid M}(U)\le r_U$ and
    $\operatorname{OSR}_{M\mid R}(V)\le r_V$.
    Define
    \begin{align}
        E_k & =\left\{\alpha:I\to\N:\sum_{a\in I}\alpha(a)=k\right\}, \notag \\
        T   & =\{\alpha\in E_k:\alpha(u)=i,\ \alpha(v)=j\}.\notag
    \end{align}
    For a word $w\in I^{\{0,\ldots,k-1\}}$, write
    $n_w(a)=|\{t:w(t)=a\}|$, and put
    \begin{align}
        C_\alpha & =\sum_{\substack{w\in I^{\{0,\ldots,k-1\}}\\n_w=\alpha}}
        H_{w(0)}\cdots H_{w(k-1)}.\notag
    \end{align}
    Every scalar family $(\lambda_\alpha)_{\alpha:I\to\N}$ satisfies
    \begin{align}
        \operatorname{OSR}_{M\mid LR}\!\left(
        \sum_{\alpha\in T}\lambda_\alpha C_\alpha\right)
        \le |T|\binom{k+i+1}{2i+1}\binom{k+j+1}{2j+1}r_U^i r_V^j.\notag
    \end{align}
    Empty slices and spaces of dimension zero are allowed, with $0^0=1$.
    These tensor hypotheses give an algebraic auxiliary to the coefficient
    selection in \cite[Lemma~4.2]{Arad2013AreaLaw} and the simultaneous
    boundary estimates in \cite[proof of Proposition~5.7]{Arad2016RigorousRG}.
    The binomial factors use the corrected count in
    \cite{gap:aklv13_interpolation_and_word_count}.
    % Physical support and variable-count identifications, low-frequency
    % selection, cut displacement and the MPO conclusion remain separate.
\end{theorem}
\begin{proof}
    \leanok
    \uses{thm:finite_monomial_interpolation,
        thm:operator_schmidt_rank_algebra,
        thm:two_marker_full_word_counts,
        thm:tripartite_polynomial_coefficient_rank}
    If $i+j>k$, the set $T$ is empty. Otherwise $i,j\le k$.
    The map $\alpha\mapsto\alpha|_O$ is injective on $T$, because the
    two omitted coordinates have the fixed values $i,j$.
    Thus the monomials $m_\alpha(z)=\prod_{a\in O}z_a^{\alpha(a)}$,
    indexed by $\alpha\in T$, are distinct.

    Put $D(z)=\sum_{a\in O}z_aH_a$ and
    $F_z(X,Y)=D(z)+H_uX+H_vY$.
    Theorem~\ref{thm:two_marker_full_word_counts} gives the same operator family
    \begin{align}
        P(z) & =[Y^j][X^i]F_z(X,Y)^k
        =\sum_{\alpha\in T}m_\alpha(z)C_\alpha.
        \label{eq:weighted_frequency_evaluation}
    \end{align}
    The local forms of $H_a$ express $D(z)$ as the sum of the three local
    operators $\sum_{a\in O}z_aA_{L,a}$,
    $\sum_{a\in O}z_aA_{M,a}$ and $\sum_{a\in O}z_aA_{R,a}$,
    each tensored with the other two identities.
    Theorem~\ref{thm:tripartite_polynomial_coefficient_rank} therefore bounds
    the middle-cut rank of every $P(z)$ by
    $R=\binom{k+i+1}{2i+1}\binom{k+j+1}{2j+1}r_U^i r_V^j$.

    Set $t=|T|$. By Theorem~\ref{thm:finite_monomial_interpolation}, choose
    points $z_0,\ldots,z_{t-1}$ and coefficients $c_{\alpha s}$ that recover
    these monomials. Applying the same recovery coefficients to
    (\ref{eq:weighted_frequency_evaluation}) gives
    \begin{align}
        \sum_{\alpha\in T}\lambda_\alpha C_\alpha
        =\sum_{s=0}^{t-1}\left(
        \sum_{\alpha\in T}\lambda_\alpha c_{\alpha s}\right)P(z_s).\notag
    \end{align}
    Scalar multiplication does not increase operator Schmidt rank.
    Subadditivity now gives the bound $tR$, using only this one set of
    $t$ evaluation points.
\end{proof}

\begin{definition}[Ordered coefficients with one distinguished interaction]
    \label{def:one_marker_word_coefficient}
    \lean{WordMultiplicity.oneMarkerCoefficient,
        WordMultiplicity.map_oneMarkerCoefficient}
    \leanok
    Let $I$ be finite and let $H_a,T$ belong to a semiring. For integers
    $m,k\ge0$ and a frequency vector $\alpha:I\to\mathbb N$, define
    \begin{align*}
        C_\alpha
        =\sum_{\substack{w\in(I\sqcup\{*\})^{\{0,\ldots,m+k-1\}} \\
                n_w(*)=k,\ n_w|_I=\alpha}}
        H_{w(0)}\cdots H_{w(m+k-1)},\qquad H_*=T.
    \end{align*}
    These coefficients commute with semiring homomorphisms.
    They are the ordered-word coefficients in the one-distinguished-interaction
    argument of \cite[Lemma~4.2]{Arad2013AreaLaw}.
\end{definition}

\begin{theorem}[Weighted frequency sums at one distinguished cut]
    \label{thm:one_marker_weighted_rank}
    \lean{Matrix.operatorSchmidtRank_sum_smul_one_marker_counts_le}
    \leanok
    \uses{def:one_marker_word_coefficient}
    Suppose that the operators $H_a$ act on a finite-dimensional complex
    bipartite space and satisfy
    $H_a=A_a\otimes\mathbb 1+\mathbb 1\otimes B_a$.
    Suppose $\operatorname{OSR}(T)\le r$, and let
    $E_m=\{\alpha:I\to\mathbb N:\sum_a\alpha(a)=m\}$.
    Every scalar family $(\lambda_\alpha)_{\alpha\in E_m}$ satisfies
    \begin{align*}
        \operatorname{OSR}\!\left(\sum_{\alpha\in E_m}\lambda_\alpha C_\alpha\right)
        \le |E_m|\binom{m+2k+1}{2k+1}r^k.
    \end{align*}
    This is the algebraic fixed-frequency estimate used in
    \cite[Lemma~4.2]{Arad2013AreaLaw}, with the corrected count from
    \cite{gap:aklv13_interpolation_and_word_count}.
    No independence of the $C_\alpha$ is required.
\end{theorem}
\begin{proof}
    \leanok
    \uses{thm:word_coefficient_expansion,thm:polynomial_commuting_gap_rank,
        thm:monomial_evaluation_weighted_rank}
    For $z\in\mathbb C^I$, put
    $A(z)=(\sum_a z_aA_a)\otimes\mathbb 1$ and
    $B(z)=\mathbb 1\otimes(\sum_a z_aB_a)$.
    These operators commute and have rank at most one. Ordered-word expansion gives
    \begin{align*}
        [X^k](A(z)+B(z)+TX)^{m+k}
        =\sum_{\alpha\in E_m}\left(\prod_a z_a^{\alpha(a)}\right)C_\alpha.
    \end{align*}
    The commuting-gap coefficient bound gives rank at most
    $\binom{m+2k+1}{2k+1}r^k$ for every $z$.
    Distinct frequency vectors give distinct monomials; weighted monomial recovery
    therefore gives the stated factor $|E_m|$.
\end{proof}

\begin{theorem}[A uniform bound on the frequency sum]
    \label{thm:aklv_frequency_sum_bound}
    \lean{Nat.aklv_frequency_sum_le}
    \leanok
    Let $q,\ell,s,r$ be natural numbers with $q\ge1$, $s\le\ell+1$, and
    $r\le\ell$. Then
    \begin{align*}
         & \sum_{u=0}^{s-1}\sum_{k=0}^{\lfloor r/s\rfloor}
        \binom{r-k+s}{s}\binom{r+k+1}{2k+1}
        q^{2k+2|u+1-\lfloor s/2\rfloor|}                                      \\
         & \hspace{2em}\le\bigl(q(\ell+1)\bigr)^{3s+6\lfloor\ell/s\rfloor+4}.
    \end{align*}
    For $s=0$ the sum is empty and natural division by zero is taken to be zero.
    These are the scalar bounds in \cite[Lemma~4.2]{Arad2013AreaLaw}, with the
    corrected count from \cite{gap:aklv13_interpolation_and_word_count}.
\end{theorem}
\begin{proof}
    \leanok
    Put $B=q(\ell+1)$.
    The inequalities $\binom{n+k}{k}\le(n+1)^k$ and
    $\binom{n}{k}\le n^k$ give
    \begin{align*}
        \binom{r-k+s}{s}    & \le B^s,      &
        \binom{r+k+1}{2k+1} & \le B^{4k+2}.
    \end{align*}
    Here $r+k+1\le(\ell+1)^2$. Since $|u+1-\lfloor s/2\rfloor|\le s$,
    the remaining factor is at most $B^{2k+2s}$.
    Each summand is therefore at most $B^{3s+6\lfloor\ell/s\rfloor+2}$.
    There are $s(\lfloor r/s\rfloor+1)\le B^2$ summands.
\end{proof}

\begin{theorem}[The effective window exponent]
    \label{thm:aklv_effective_width_bound}
    \lean{Nat.aklv_effective_width_bounds}
    \leanok
    Let $\ell,s\ge1$ be integers and put
    $w=\min\{s,2\lfloor\sqrt\ell\rfloor\}$. Then $1\le w\le\ell+1$ and
    \begin{align*}
        3w+6\lfloor\ell/w\rfloor+5
        \le16\bigl(\lfloor\ell/s\rfloor+\lfloor\sqrt\ell\rfloor+1\bigr).
    \end{align*}
    This is the numerical shorter-window estimate in
    \cite[Lemma~4.2]{Arad2013AreaLaw}.
\end{theorem}
\begin{proof}
    \leanok
    Set $n=\lfloor\sqrt\ell\rfloor\ge1$.
    The inequalities $n^2\le\ell\le n^2+2n$ imply $2n\le\ell+1$.
    If $w=s\le2n$, substitute this width bound.
    Otherwise $w=2n$ and
    $\ell\le n^2+2n\le2n(n+1)$ gives
    $\lfloor\ell/w\rfloor\le n+1$. Substitution gives the assertion.
\end{proof}
